% Глава 1. Анализ предметной области и постановка задачи.
% Источники: knowledge-base/01-problem-and-audience.md,
%            knowledge-base/02-competitors.md,
%            knowledge-base/03-history-and-milestones.md.

\clearpage
\section{Анализ предметной области и постановка задачи}

\subsection{Учебное расписание как объект цифровизации}

% TODO: контент
Учебное расписание является центральным операционным артефактом образовательного процесса: на его основе планируется работа всех участников --- от студентов и преподавателей до администрации учебного заведения. При этом в значительной части российских вузов работа с расписанием остаётся ручной: составление ведётся в Excel-таблицах, публикация --- в виде PDF-файлов на сайтах вузов, а оперативные изменения распространяются через старост и неофициальные мессенджер-каналы. Такая модель создаёт системные потери на каждом этапе: рост числа ресурсных коллизий при ручном составлении, задержка в доведении изменений до конечных пользователей и отсутствие обратной связи о реальном использовании опубликованных данных.

\subsection{Целевые роли и сценарии использования}

\subsubsection{Студент}
% TODO: контент
Студент --- ежедневный потребитель расписания. Ключевые требования: мгновенный доступ к расписанию своей группы и интересующих преподавателей, корректная работа в условиях нестабильного интернет-соединения, видимость оперативных изменений (переносы, замены, изменения аудиторий), возможность персонализации (тема, акцентные цвета, заметки к парам). В существующих каналах распространения расписания эти требования не удовлетворены: сайты вузов плохо адаптированы под мобильные устройства, не работают офлайн и не поддерживают пользовательские заметки. Приложения-парсеры закрывают часть UX-потребностей, но критически зависят от стабильности сайта-источника и монетизируются за счёт рекламы, что снижает качество пользовательского опыта.

\subsubsection{Преподаватель}
% TODO: контент
Сценарии преподавателя в значительной мере совпадают со студенческими. Дополнительные требования: удобный обзор собственной учебной нагрузки, быстрое уведомление группы о переносе пары, а в перспективе --- инструменты отметки посещаемости. Существующие решения, ориентированные на студента, для преподавателя малопригодны.

\subsubsection{Методист и администрация учебного заведения}
% TODO: контент
На стороне учебного заведения работа с расписанием ведётся методистами и диспетчерами учебного отдела. Типовая задача --- оперативное изменение уже опубликованной сетки занятий --- требует одновременной координации по нескольким измерениям: загруженность преподавателей, аудиторного фонда, учебных групп и аудиторных требований дисциплин. Ручное ведение в Excel не обеспечивает встроенной проверки коллизий, а доставка обновлений студентам становится отдельной операцией. Прямые конкуренты на этом сегменте рынка (см. п.~\ref{sec:direct-competitors}) предлагают on-premise решения, требующие собственной инфраструктуры и квалифицированного сопровождения.

\subsection{Обзор существующих решений}

\subsubsection{Прямые конкуренты}
\label{sec:direct-competitors}
% TODO: контент. Подробные таблицы — knowledge-base/02-competitors.md.
В рамках работы рассмотрены три прямых конкурента, действующих на российском рынке систем управления учебным расписанием: <<Галактика РУЗ>>, <<ММИС~Лаб>> и <<1С:Университет>>. Все три решения поставляются по on-premise модели, не имеют нативной мобильной составляющей и интегрированы в более широкие ERP-экосистемы соответствующих вендоров. Подробное сопоставление по параметрам поставки, мобильной части, интеграции, стоимости и целевой аудитории представлено далее в работе.

% TODO: вставить таблицу прямых конкурентов из knowledge-base/02-competitors.md

\subsubsection{Непрямые конкуренты --- приложения-парсеры}
% TODO: контент
Отдельную нишу занимают мобильные приложения, не управляющие расписанием, а агрегирующие его путём парсинга публичных сайтов учебных заведений (наиболее показательный пример --- приложение <<Кампус>>). Они подтверждают наличие массового спроса на мобильный доступ к расписанию, однако имеют ряд системных ограничений: нестабильность парсинга при изменении сайта-источника, отсутствие реального времени, рекламная монетизация, отсутствие институциональной легитимности и односторонняя функциональность (только просмотр).

\subsubsection{Сводное сравнение}
% TODO: вставить сводную таблицу сравнения из knowledge-base/02-competitors.md
Сводное сравнение по ключевым параметрам выбора решения учебным заведением показывает, что ни один из прямых конкурентов одновременно не поддерживает все характеристики, формирующие современное требование рынка: SaaS-модель поставки, нативные мобильные приложения, современный UI/UX, открытый API, независимость от сторонних ERP-экосистем и низкий порог входа по стоимости.

\subsection{Выводы по анализу и требования к клиентской части}

\subsubsection{Функциональные требования}
% TODO: контент
По итогам анализа сформулированы функциональные требования к клиентской части платформы. Мобильные приложения должны обеспечивать просмотр расписания с возможностью офлайн-доступа, персонализацию (тема, акцентные цвета, иконка приложения, язык), работу с избранными сущностями (группы, преподаватели, аудитории), возможность оставлять заметки к парам и виджеты на рабочем столе. Веб-приложение должно поддерживать публичную часть (лендинг, веб-просмотр расписания) и закрытый сервис администрирования с управлением справочниками учебного заведения, сотрудниками и составлением расписания.

\subsubsection{Нефункциональные требования}
% TODO: контент
Нефункциональные требования включают: соответствие нативным дизайн-гайдлайнам каждой платформы (Apple HIG, Material Design 3 Expressive, shadcn/ui), офлайн-устойчивость мобильных клиентов, время отклика интерфейса в пределах требований платформ к плавности интерфейса, мультиязычность (русский и английский), доступность с устройств на iOS~18 и выше, Android (минимальный SDK уточняется в~гл.~3) и современных веб-браузеров.
